\documentclass[10pt,a4paper]{amsart}
\usepackage{amsmath,amssymb,amsthm,graphicx}
% Retained minimal preamble for the Stacks Project chapter varieties.tex.
% Verbatim excerpt of the chapter preamble: Xy-pic setup (native preamble lines 14-18), the native theorem
% declarations with the mathematical macros (native preamble lines 154-217), and the project own no-op
% definitions of the reference/slogan/history display environments (native wrapper.tex lines 1-3) that the
% counted statement uses for its slogan. Copied mechanically; nothing added, removed or redefined.
% For commutative diagrams we use Xy-pic
\usepackage[all]{xy}

% We use 2cell for 2-commutative diagrams.
\xyoption{2cell}
% Theorem environments.
%
\theoremstyle{plain}
\newtheorem{theorem}[subsection]{Theorem}
\newtheorem{proposition}[subsection]{Proposition}
\newtheorem{lemma}[subsection]{Lemma}

\theoremstyle{definition}
\newtheorem{definition}[subsection]{Definition}
\newtheorem{example}[subsection]{Example}
\newtheorem{exercise}[subsection]{Exercise}
\newtheorem{situation}[subsection]{Situation}

\theoremstyle{remark}
\newtheorem{remark}[subsection]{Remark}
\newtheorem{remarks}[subsection]{Remarks}

\numberwithin{equation}{subsection}

% Macros
%
\def\lim{\mathop{\mathrm{lim}}\nolimits}
\def\colim{\mathop{\mathrm{colim}}\nolimits}
\def\Spec{\mathop{\mathrm{Spec}}}
\def\Hom{\mathop{\mathrm{Hom}}\nolimits}
\def\Ext{\mathop{\mathrm{Ext}}\nolimits}
\def\SheafHom{\mathop{\mathcal{H}\!\mathit{om}}\nolimits}
\def\SheafExt{\mathop{\mathcal{E}\!\mathit{xt}}\nolimits}
\def\Sch{\mathit{Sch}}
\def\Mor{\mathop{\mathrm{Mor}}\nolimits}
\def\Ob{\mathop{\mathrm{Ob}}\nolimits}
\def\Sh{\mathop{\mathit{Sh}}\nolimits}
\def\NL{\mathop{N\!L}\nolimits}
\def\CH{\mathop{\mathrm{CH}}\nolimits}
\def\proetale{{pro\text{-}\acute{e}tale}}
\def\etale{{\acute{e}tale}}
\def\QCoh{\mathit{QCoh}}
\def\Ker{\mathop{\mathrm{Ker}}}
\def\Im{\mathop{\mathrm{Im}}}
\def\Coker{\mathop{\mathrm{Coker}}}
\def\Coim{\mathop{\mathrm{Coim}}}

% Boxtimes
%
\DeclareMathSymbol{\boxtimes}{\mathbin}{AMSa}{"02}

%
% Macros for moduli stacks/spaces
%
\def\QCohstack{\mathcal{QC}\!\mathit{oh}}
\def\Cohstack{\mathcal{C}\!\mathit{oh}}
\def\Spacesstack{\mathcal{S}\!\mathit{paces}}
\def\Quotfunctor{\mathrm{Quot}}
\def\Hilbfunctor{\mathrm{Hilb}}
\def\Curvesstack{\mathcal{C}\!\mathit{urves}}
\def\Polarizedstack{\mathcal{P}\!\mathit{olarized}}
\def\Complexesstack{\mathcal{C}\!\mathit{omplexes}}
% \Pic is the operator that assigns to X its picard group, usage \Pic(X)
% \Picardstack_{X/B} denotes the Picard stack of X over B
% \Picardfunctor_{X/B} denotes the Picard functor of X over B
\def\Pic{\mathop{\mathrm{Pic}}\nolimits}
\def\Picardstack{\mathcal{P}\!\mathit{ic}}
\def\Picardfunctor{\mathrm{Pic}}
\def\Deformationcategory{\mathcal{D}\!\mathit{ef}}
\newenvironment{reference}{}{}
\newenvironment{slogan}{}{}
\newenvironment{history}{}{}
\usepackage{cleveref}
\crefname{originallemma}{original lemma}{original lemmas}
\Crefname{originallemma}{Original lemma}{Original lemmas}
\makeatletter
\newlabel{lemma-kunneth-general}{{0BEF}{source}{}{sourcepointer-1}{}}
\newlabel{lemma-kunneth-general@cref}{{[originallemma][0][]0BEF}{[1][0][]source}{}{}{}}
\makeatother
\crefname{originalreference}{original reference}{original references}
\Crefname{originalreference}{Original reference}{Original references}
\makeatletter
\newlabel{coherent-lemma-quasi-coherent-affine-cohomology-zero}{{01XB}{source}{}{sourcepointer-2}{}}
\newlabel{coherent-lemma-quasi-coherent-affine-cohomology-zero@cref}{{[originalreference][0][]01XB}{[1][0][]source}{}{}{}}
\makeatother
\crefname{originalreference}{original reference}{original references}
\Crefname{originalreference}{Original reference}{Original references}
\makeatletter
\newlabel{cohomology-lemma-functoriality-cech}{{01FD}{source}{}{sourcepointer-3}{}}
\newlabel{cohomology-lemma-functoriality-cech@cref}{{[originalreference][0][]01FD}{[1][0][]source}{}{}{}}
\makeatother
\crefname{originalreference}{original reference}{original references}
\Crefname{originalreference}{Original reference}{Original references}
\makeatletter
\newlabel{cohomology-equation-needs-signs}{{07MB}{source}{}{sourcepointer-4}{}}
\newlabel{cohomology-equation-needs-signs@cref}{{[originalreference][0][]07MB}{[1][0][]source}{}{}{}}
\makeatother
\crefname{originalreference}{original reference}{original references}
\Crefname{originalreference}{Original reference}{Original references}
\makeatletter
\newlabel{coherent-lemma-affine-diagonal}{{0BDX}{source}{}{sourcepointer-5}{}}
\newlabel{coherent-lemma-affine-diagonal@cref}{{[originalreference][0][]0BDX}{[1][0][]source}{}{}{}}
\makeatother
\crefname{originalreference}{original reference}{original references}
\Crefname{originalreference}{Original reference}{Original references}
\makeatletter
\newlabel{coherent-lemma-cech-cohomology-quasi-coherent}{{01XD}{source}{}{sourcepointer-6}{}}
\newlabel{coherent-lemma-cech-cohomology-quasi-coherent@cref}{{[originalreference][0][]01XD}{[1][0][]source}{}{}{}}
\makeatother
\crefname{originalreference}{original reference}{original references}
\Crefname{originalreference}{Original reference}{Original references}
\makeatletter
\newlabel{coherent-lemma-vanishing-nr-affines}{{01XI}{source}{}{sourcepointer-7}{}}
\newlabel{coherent-lemma-vanishing-nr-affines@cref}{{[originalreference][0][]01XI}{[1][0][]source}{}{}{}}
\makeatother
\crefname{originalreference}{original reference}{original references}
\Crefname{originalreference}{Original reference}{Original references}
\makeatletter
\newlabel{coherent-lemma-affine-diagonal-universal-delta-functor}{{0BDY}{source}{}{sourcepointer-8}{}}
\newlabel{coherent-lemma-affine-diagonal-universal-delta-functor@cref}{{[originalreference][0][]0BDY}{[1][0][]source}{}{}{}}
\makeatother
\crefname{originalreference}{original reference}{original references}
\Crefname{originalreference}{Original reference}{Original references}
\makeatletter
\newlabel{cohomology-lemma-quasi-separated-cohomology-colimit}{{01FF}{source}{}{sourcepointer-9}{}}
\newlabel{cohomology-lemma-quasi-separated-cohomology-colimit@cref}{{[originalreference][0][]01FF}{[1][0][]source}{}{}{}}
\makeatother
\crefname{originalreference}{original reference}{original references}
\Crefname{originalreference}{Original reference}{Original references}
\makeatletter
\newlabel{cohomology-lemma-apply-Leray}{{01F4}{source}{}{sourcepointer-10}{}}
\newlabel{cohomology-lemma-apply-Leray@cref}{{[originalreference][0][]01F4}{[1][0][]source}{}{}{}}
\makeatother
\title{The Kunneth formula for the cohomology of a product of quasi-compact schemes with affine diagonal}
\author{The Stacks Project Authors}
\date{stacks-project:varieties.tex:local-snapshot-20261011; GNU FDL 1.2 or later}
\begin{document}
\maketitle
\noindent\textit{Editorial note.} All delivered mathematical text is the author's own, copied verbatim from the pinned Stacks Project chapter file varieties.tex; no statement or proof step was rewritten, reordered or removed.\par\medskip
\noindent\textit{Editorial note.} The counted span cites twelve cross-chapter labels, each recorded in assembly.external\_references with its canonical tag as a navigation pointer; none replaces a load-bearing step, which are all inside the span.\par\medskip
\noindent\textit{Editorial note.} The proof cites no bibliography item inside the counted span.\par\medskip
\noindent\textit{Allowed standard prerequisites (editorial).}\begin{itemize}
\item Vanishing of higher cohomology of quasi-coherent modules on affine schemes. Stacks Project, Cohomology of Schemes, Lemma coherent-lemma-quasi-coherent-affine-cohomology-zero (tag 01XB). Supplies the affine case of the theorem.
\item Functoriality of Cech cohomology and Cech-to-cohomology comparison. Stacks Project, Cohomology, Lemma cohomology-lemma-functoriality-cech (tag 01EO). Produces the maps between Cech complexes over the three coverings used in the proof.
\end{itemize}\medskip
\ifdefined\hypertarget\hypertarget{sourcepointer-1}{}\fi
\noindent\textit{Source reference.} Original lemma 0BEF: sources/stacks-chow/tags.txt, source lines 14348--14348. Cross-chapter standard prerequisite cited by the author inside the counted span; it is used as a pointer and its own proof is not part of this item.\par\medskip
\ifdefined\hypertarget\hypertarget{sourcepointer-2}{}\fi
\noindent\textit{Source reference.} Original reference 01XB: sources/stacks-chow/tags.txt, source lines 2512--2512. Cross-chapter standard prerequisite cited by the author inside the counted span; it is used as a pointer and its own proof is not part of this item.\par\medskip
\ifdefined\hypertarget\hypertarget{sourcepointer-3}{}\fi
\noindent\textit{Source reference.} Original reference 01FD: sources/stacks-chow/tags.txt, source lines 1866--1866. Cross-chapter standard prerequisite cited by the author inside the counted span; it is used as a pointer and its own proof is not part of this item.\par\medskip
\ifdefined\hypertarget\hypertarget{sourcepointer-4}{}\fi
\noindent\textit{Source reference.} Original reference 07MB: sources/stacks-chow/tags.txt, source lines 9721--9721. Cross-chapter standard prerequisite cited by the author inside the counted span; it is used as a pointer and its own proof is not part of this item.\par\medskip
\ifdefined\hypertarget\hypertarget{sourcepointer-5}{}\fi
\noindent\textit{Source reference.} Original reference 0BDX: sources/stacks-chow/tags.txt, source lines 14330--14330. Cross-chapter standard prerequisite cited by the author inside the counted span; it is used as a pointer and its own proof is not part of this item.\par\medskip
\ifdefined\hypertarget\hypertarget{sourcepointer-6}{}\fi
\noindent\textit{Source reference.} Original reference 01XD: sources/stacks-chow/tags.txt, source lines 2514--2514. Cross-chapter standard prerequisite cited by the author inside the counted span; it is used as a pointer and its own proof is not part of this item.\par\medskip
\ifdefined\hypertarget\hypertarget{sourcepointer-7}{}\fi
\noindent\textit{Source reference.} Original reference 01XI: sources/stacks-chow/tags.txt, source lines 2519--2519. Cross-chapter standard prerequisite cited by the author inside the counted span; it is used as a pointer and its own proof is not part of this item.\par\medskip
\ifdefined\hypertarget\hypertarget{sourcepointer-8}{}\fi
\noindent\textit{Source reference.} Original reference 0BDY: sources/stacks-chow/tags.txt, source lines 14331--14331. Cross-chapter standard prerequisite cited by the author inside the counted span; it is used as a pointer and its own proof is not part of this item.\par\medskip
\ifdefined\hypertarget\hypertarget{sourcepointer-9}{}\fi
\noindent\textit{Source reference.} Original reference 01FF: sources/stacks-chow/tags.txt, source lines 1868--1868. Cross-chapter standard prerequisite cited by the author inside the counted span; it is used as a pointer and its own proof is not part of this item.\par\medskip
\ifdefined\hypertarget\hypertarget{sourcepointer-10}{}\fi
\noindent\textit{Source reference.} Original reference 01F4: sources/stacks-chow/tags.txt, source lines 1857--1857. Cross-chapter standard prerequisite cited by the author inside the counted span; it is used as a pointer and its own proof is not part of this item.\par\medskip
\begin{lemma}
\label{lemma-kunneth}
Let $k$ be a field. Let $X$ and $Y$ be schemes over $k$ and
let $\mathcal{F}$, resp.\ $\mathcal{G}$ be a quasi-coherent
$\mathcal{O}_X$-module, resp.\ $\mathcal{O}_Y$-module.
Then we have a canonical isomorphism
$$
H^n(X \times_{\Spec(k)} Y, \text{pr}_1^*\mathcal{F}
\otimes_{\mathcal{O}_{X \times_{\Spec(k)} Y}} \text{pr}_2^*\mathcal{G}) =
\bigoplus\nolimits_{p + q = n}
H^p(X, \mathcal{F}) \otimes_k H^q(Y, \mathcal{G})
$$
provided $X$ and $Y$ are quasi-compact and have affine
diagonal\footnote{The case where $X$ and $Y$ are quasi-separated
will be discussed in Lemma \ref{lemma-kunneth-general} below.}
(for example if $X$ and $Y$ are separated).
\end{lemma}

\begin{proof}
In this proof unadorned products and tensor products are over $k$. As maps
$$
H^p(X, \mathcal{F}) \otimes H^q(Y, \mathcal{G})
\longrightarrow
H^n(X \times Y, \text{pr}_1^*\mathcal{F}
\otimes_{\mathcal{O}_{X \times Y}} \text{pr}_2^*\mathcal{G})
$$
we use functoriality of cohomology to get maps
$H^p(X, \mathcal{F}) \to H^p(X \times Y, \text{pr}_1^*\mathcal{F})$ and
$H^p(Y, \mathcal{G}) \to H^p(X \times Y, \text{pr}_2^*\mathcal{G})$
and then we use the cup product
$$
\cup :
H^p(X \times Y, \text{pr}_1^*\mathcal{F})
\otimes H^q(X \times Y, \text{pr}_2^*\mathcal{G})
\longrightarrow
H^n(X \times Y, \text{pr}_1^*\mathcal{F}
\otimes_{\mathcal{O}_{X \times Y}} \text{pr}_2^*\mathcal{G})
$$
The result is true when $X$ and $Y$ are affine by the vanishing of
higher cohomology groups on affines (Cohomology of Schemes, Lemma
\ref{coherent-lemma-quasi-coherent-affine-cohomology-zero})
and the definitions (of pullbacks of quasi-coherent modules and tensor
products of quasi-coherent modules).

\medskip\noindent
Choose finite affine open coverings
$\mathcal{U} : X = \bigcup_{i \in I} U_i$ and
$\mathcal{V} : Y = \bigcup_{j \in J} V_j$.
This determines an affine open covering
$\mathcal{W} : X \times Y = \bigcup_{(i, j) \in I \times J} U_i \times V_j$.
Note that $\mathcal{W}$ is a refinement of
$\text{pr}_1^{-1}\mathcal{U}$ and of $\text{pr}_2^{-1}\mathcal{V}$.
Thus by Cohomology, Lemma \ref{cohomology-lemma-functoriality-cech}
we obtain maps
$$
\check{\mathcal{C}}^\bullet(\mathcal{U}, \mathcal{F}) \to
\check{\mathcal{C}}^\bullet(\mathcal{W}, \text{pr}_1^*\mathcal{F})
\quad\text{and}\quad
\check{\mathcal{C}}^\bullet(\mathcal{V}, \mathcal{G}) \to
\check{\mathcal{C}}^\bullet(\mathcal{W}, \text{pr}_2^*\mathcal{G})
$$
compatible with pullback maps on cohomology. In Cohomology, Equation
(\ref{cohomology-equation-needs-signs})
we have constructed a map of complexes
$$
\text{Tot}(
\check{\mathcal{C}}^\bullet(\mathcal{W}, \text{pr}_1^*\mathcal{F})
\otimes
\check{\mathcal{C}}^\bullet(\mathcal{W}, \text{pr}_2^*\mathcal{G}))
\longrightarrow
\check{\mathcal{C}}^\bullet(\mathcal{W},
\text{pr}_1^*\mathcal{F} \otimes_{\mathcal{O}_{X \times Y}}
\text{pr}_2^*\mathcal{G})
$$
defining the cup product on cohomology. Combining the above we
obtain a map of complexes
\begin{equation}
\label{equation-kunneth-on-cech}
\text{Tot}(
\check{\mathcal{C}}^\bullet(\mathcal{U}, \mathcal{F})
\otimes
\check{\mathcal{C}}^\bullet(\mathcal{V}, \mathcal{G}))
\longrightarrow
\check{\mathcal{C}}^\bullet(\mathcal{W},
\text{pr}_1^*\mathcal{F} \otimes_{\mathcal{O}_{X \times Y}}
\text{pr}_2^*\mathcal{G})
\end{equation}
We warn the reader that this map is not an isomorphism of
complexes. Recall that we may compute the cohomologies of our
quasi-coherent sheaves using our coverings
(Cohomology of Schemes, Lemmas
\ref{coherent-lemma-affine-diagonal} and
\ref{coherent-lemma-cech-cohomology-quasi-coherent}).
Thus on cohomology (\ref{equation-kunneth-on-cech}) reproduces the map of
the lemma.

\medskip\noindent
Consider a short exact
sequence $0 \to \mathcal{F} \to \mathcal{F}' \to \mathcal{F}'' \to 0$
of quasi-coherent modules. Since the construction of
(\ref{equation-kunneth-on-cech}) is functorial in $\mathcal{F}$ and since the
formation of the relevant {\v C}ech complexes is exact in
the variable $\mathcal{F}$ (because we are taking sections over
affine opens) we find a map between short exact sequence of
complexes
$$
\xymatrix{
\text{Tot}(
\check{\mathcal{C}}^\bullet(\mathcal{U}, \mathcal{F})
\otimes
\check{\mathcal{C}}^\bullet(\mathcal{V}, \mathcal{G})) \ar[r] \ar[d] &
\text{Tot}(
\check{\mathcal{C}}^\bullet(\mathcal{U}, \mathcal{F}')
\otimes
\check{\mathcal{C}}^\bullet(\mathcal{V}, \mathcal{G})) \ar[r] \ar[d] &
\text{Tot}(
\check{\mathcal{C}}^\bullet(\mathcal{U}, \mathcal{F}'')
\otimes
\check{\mathcal{C}}^\bullet(\mathcal{V}, \mathcal{G})) \ar[d] \\
\check{\mathcal{C}}^\bullet(\mathcal{W},
\text{pr}_1^*\mathcal{F} \otimes_{\mathcal{O}_{X \times Y}}
\text{pr}_2^*\mathcal{G}) \ar[r] &
\check{\mathcal{C}}^\bullet(\mathcal{W},
\text{pr}_1^*\mathcal{F}' \otimes_{\mathcal{O}_{X \times Y}}
\text{pr}_2^*\mathcal{G}) \ar[r] &
\check{\mathcal{C}}^\bullet(\mathcal{W},
\text{pr}_1^*\mathcal{F}'' \otimes_{\mathcal{O}_{X \times Y}}
\text{pr}_2^*\mathcal{G})
}
$$
(we have dropped the outer zeros).
Looking at long exact cohomology sequences we find that if the result of
the lemma holds for $2$-out-of-$3$ of
$\mathcal{F}, \mathcal{F}', \mathcal{F}''$, then it holds for
the third.

\medskip\noindent
Observe that $X$ has finite cohomological dimension for
quasi-coherent modules, see Cohomology of Schemes, Lemma
\ref{coherent-lemma-vanishing-nr-affines}.
Using induction on
$d(\mathcal{F}) = \max \{d \mid H^d(X, \mathcal{F}) \not = 0\}$
we will reduce to the case $d(\mathcal{F}) = 0$.
Assume $d(\mathcal{F}) > 0$.
By Cohomology of Schemes, Lemma
\ref{coherent-lemma-affine-diagonal-universal-delta-functor}
we have seen that
there exists an embedding $\mathcal{F} \to \mathcal{F}'$
such that $H^p(X, \mathcal{F}') = 0$ for all $p \geq 1$.
Setting $\mathcal{F}'' = \Coker(\mathcal{F} \to \mathcal{F}')$
we see that $d(\mathcal{F}'') < d(\mathcal{F})$.
Then we can apply the result from the previous paragraph
to see that it suffices to prove the lemma for $\mathcal{F}'$
and $\mathcal{F}''$ thereby proving the induction step.

\medskip\noindent
Arguing in the same fashion for $\mathcal{G}$ we find that we
may assume that both $\mathcal{F}$ and $\mathcal{G}$
have nonzero cohomology only in degree $0$.
Let $V \subset Y$ be an affine open. Consider the affine open covering
$\mathcal{U}_V : X \times V = \bigcup_{i \in I} U_i \times V$.
It is immediate that
$$
\check{\mathcal{C}}^\bullet(\mathcal{U}, \mathcal{F}) \otimes \mathcal{G}(V) =
\check{\mathcal{C}}^\bullet(\mathcal{U}_V,
\text{pr}_1^*\mathcal{F} \otimes_{\mathcal{O}_{X \times Y}}
\text{pr}_2^*\mathcal{G})
$$
(equality of complexes). We conclude that
$$
R\text{pr}_{2, *}(\text{pr}_1^*\mathcal{F} \otimes_{\mathcal{O}_{X \times Y}}
\text{pr}_2^*\mathcal{G})
\cong
\Gamma(X, \mathcal{F}) \otimes_k \mathcal{G} \cong
\bigoplus\nolimits_{\alpha \in A} \mathcal{G}
$$
on $Y$. Here $A$ is a basis for the $k$-vector space $\Gamma(X, \mathcal{F})$.
Cohomology on $Y$ commutes with direct sums
(Cohomology, Lemma \ref{cohomology-lemma-quasi-separated-cohomology-colimit}).
Using the Leray spectral sequence for $\text{pr}_2$
(via Cohomology, Lemma \ref{cohomology-lemma-apply-Leray})
we conclude that
$H^n(X \times Y, \text{pr}_1^*\mathcal{F} \otimes_{\mathcal{O}_{X \times Y}}
\text{pr}_2^*\mathcal{G})$
is zero for $n > 0$ and isomorphic to
$H^0(X, \mathcal{F}) \otimes H^0(Y, \mathcal{G})$ for $n = 0$.
This finishes the proof (except that we should check that the
isomorphism is indeed given by cup product in degree $0$; we omit
the verification).
\end{proof}

\end{document}
